\section{Conclusions}

Embodied action and augmented reality are separated from the ARC-ADHD core intervention because they create distinct causal, technical, and safety questions. The proposed study distinguishes stationary interaction, embodied movement without AR, and embodied AR so that movement effects are not mislabeled as AR effects.

The central requirement is conservative: adding embodiment or AR is useful only if the intended contextual or movement effect occurs without an unacceptable loss of cognitive control or safety. A later longitudinal study is required to determine whether context variation improves independent transfer.

This companion design allows the main ARC-ADHD framework to remain minimal while preserving a rigorous path for testing whether moving the same regulation rule from screen to physical space provides added value.
